
\uline{régulier}, \uline{si}, \uline{localement sur} $X$, \uline{il existe un} $\underline{A}$-\uline{Module libre de type fini} $\underline{E}$, \uline{et un homomorphisme régulier surjectif} : $\underline{E}\longrightarrow \underline{J}$.

\uline{Soit} $i:Y\longrightarrow X$ \uline{une immersion de préschémas}, \uline{et soit} $U$ \uline{un ouvert de} $X$ \uline{tel que} $i(Y)\subset U$ \uline{et que} $i$ \uline{soit une immersion fermée de} $Y$ \uline{dans} $U$. \uline{On dit que} $i$ \uline{est régulière si l'idéal de} $\underline{O}_U$ \uline{définissant} $i(Y)$ \uline{est régulier} (\uline{condition qui ne dépend pas de} $U$). \uline{Si} $y\in Y$, \uline{on dit que} $i$ \uline{est régulière en} $y$ \uline{s'il existe un voisinage} $V$ \uline{de} $i(y)$ \uline{tel que l'immersion} $i^{-1}(V)\longrightarrow V$ \uline{soit régulière}. \uline{Pour qu'une immersion soit régulière}, \uline{il est donc nécéssaire et suffisant qu'elle le soit au voisinage de chaque point}.

\emph{Cette définition est en conflit avec} EGA IV 16.9.2, \emph{et remplace cette dernière}. Toutefois, dans le cas noethérien, les deux définitions sont équivalentes (EGA IV 19.5.1). Une immersion qui peut localement être définie par une suite régulière est régulière; une immersion régulière est quasi-régulière (1.3).

\textbf{1.4.1.} Soient $X$ un schéma, $\underline{J}$ un Idéal de type fini de $\underline{O}_X$, $x$ un point de $\Supp(\underline{O}_X/\underline{J})$. Nous dirons qu'une famille de sections $f_1,\ldots,f_d$ de $\underline{J}$ sur un voisinage de $x$ forme un \uline{système minimum de générateurs de} $\underline{J}$ \uline{au voisinage de} $x$ si les $(f_i)_x$ sont une famille de générateurs de $\underline{J}_x$ de nombre minimum, i.e. leurs images dans $\underline{J}_x\otimes k(x)$ forment une base de cet espace vectoriel sur $k(x)$. On notera que les $f_i$ engendrent $\underline{J}$ sur un voisinage de $x$.

\textbf{1.4.2.} Dire qu'une immersion est régulière, c'est encore dire qu'elle est définie dans un certain ouvert de $X$ par un idéal de type fini $\underline{J}$, et que pour tout point $x\in\Supp(\underline{O}_X/\underline{J})$, et tout système minimum de générateurs de $\underline{J}$ au voisinage de $x$, le complexe de Koszul associé à cette suite est acyclique au voisinage de $x$: soit en effet $f_1,\ldots,f_d$ un système minimum de
